\documentclass[11pt, a4paper]{article}

% Gemeinsame Style-Definitionen (TEXINPUTS muss templates/ enthalten — siehe scripts/build_tex.py)
% bfw-style.tex — gemeinsame Style-Definitionen für alle Handouts/Aufgabenhefte
% Voraussetzung: Kompilation mit xelatex (wegen fontspec)

\usepackage[ngerman]{babel}
% Babel-ngerman macht " zu einem aktiven Shorthand (z.B. für "a -> ä). In unseren
% Texten verwenden wir " als normales Zeichen — daher Shorthand komplett ausschalten.
\shorthandoff{"}
\usepackage{geometry}
\geometry{a4paper, top=2.4cm, bottom=2.4cm, left=2.3cm, right=2.3cm,
          headheight=15pt, headsep=0.7cm, footskip=1.1cm}

\usepackage{fontspec}
% Fallback-Kette: Source Sans Pro -> Calibri -> Segoe UI -> Latin Modern Sans
% (Latin Modern ist immer mit MiKTeX vorhanden)
\IfFontExistsTF{Source Sans Pro}{%
  \setmainfont{Source Sans Pro}[Ligatures=TeX]%
}{\IfFontExistsTF{Calibri}{%
  \setmainfont{Calibri}[Ligatures=TeX]%
}{\IfFontExistsTF{Segoe UI}{%
  \setmainfont{Segoe UI}[Ligatures=TeX]%
}{%
  \setmainfont{Latin Modern Sans}[Ligatures=TeX]%
}}}
\IfFontExistsTF{Source Code Pro}{%
  \setmonofont{Source Code Pro}[Scale=0.92]%
}{\IfFontExistsTF{Consolas}{%
  \setmonofont{Consolas}[Scale=0.92]%
}{%
  \setmonofont{Latin Modern Mono}[Scale=0.92]%
}}

% Gesperrte Versalien für Labels/Titelbalken (Kapitälchen-Ersatz, funktioniert
% mit jeder OpenType-Schrift — Latin Modern Sans hat keine echten Kapitälchen)
\newcommand{\bfwlabelfont}{\footnotesize\bfseries\addfontfeatures{LetterSpace=8.0}}

\usepackage{xcolor}
% Farbpalette: hell, freundlich, modern
\definecolor{bfwPrimary}{HTML}{0EA5A4}    % Petrol/Türkis - Primär
\definecolor{bfwPrimaryDark}{HTML}{0F766E} % dunkler Petrol für Headings
\definecolor{bfwAccent}{HTML}{F59E0B}     % Amber - Akzent
\definecolor{bfwSuccess}{HTML}{10B981}    % Grün - Erfolg / Lösung
\definecolor{bfwWarn}{HTML}{F97316}       % Orange - Achtung
\definecolor{bfwInk}{HTML}{0F172A}        % fast Schwarz
\definecolor{bfwInkSoft}{HTML}{475569}    % gedämpftes Grau
\definecolor{bfwBgSoft}{HTML}{F8FAFC}     % off-white
\definecolor{bfwBgTeal}{HTML}{ECFEFF}     % helles Türkis
\definecolor{bfwBgAmber}{HTML}{FEF3C7}    % helles Gelb
\definecolor{bfwBgRose}{HTML}{FFE4E6}     % helles Rosé für Eselsbrücken
\definecolor{bfwTabHead}{HTML}{E4F3F2}    % zarte Kopfzeilen-Hinterlegung für Tabellen

\usepackage[colorlinks=true, linkcolor=bfwPrimaryDark, urlcolor=bfwPrimaryDark]{hyperref}
\usepackage{lastpage}
\usepackage{enumitem}
\setlist[itemize]{leftmargin=*, itemsep=2pt, topsep=2pt}
\setlist[enumerate]{leftmargin=*, itemsep=2pt, topsep=2pt}

\usepackage[most]{tcolorbox}
\usepackage{booktabs}
\usepackage{colortbl}
\usepackage{tikz}
\usetikzlibrary{decorations.pathmorphing, positioning, shapes.geometric, arrows.meta}

% ---------------------------------------------------------------------------
% Einheitliches Tabellen-Design
%   - etwas mehr Zeilenluft, dezente graue Linien (wirkt auch mit \hline ruhiger)
%   - Kopfzeile einer Tabelle bei Bedarf zart hinterlegen: \rowcolor{bfwTabHead}
% ---------------------------------------------------------------------------
\renewcommand{\arraystretch}{1.25}
\arrayrulecolor{bfwInkSoft!50}
\setlength{\heavyrulewidth}{1pt}
\setlength{\lightrulewidth}{0.5pt}

% ---------------------------------------------------------------------------
% Boxen — klare farbige Titelbalken mit gesperrten Versal-Labels
% (Titel bleibt per Option überschreibbar: \begin{merksatz}[title={...}])
% ---------------------------------------------------------------------------
\tcbset{bfwbox/.style={%
  enhanced, breakable,
  boxrule=0.5pt, arc=2.5pt,
  left=10pt, right=10pt, top=8pt, bottom=8pt,
  toptitle=4pt, bottomtitle=4pt,
  titlerule=0pt,
  fonttitle=\bfwlabelfont,
  coltitle=white
}}

% Merkkästchen — wichtige Definition / Kernpunkt
\newtcolorbox{merksatz}[1][]{%
  bfwbox,
  colback=bfwBgTeal,
  colframe=bfwPrimaryDark,
  colbacktitle=bfwPrimaryDark,
  title={MERKE},
  #1
}

% Eselsbrücke — Mnemoniks
\newtcolorbox{eselsbruecke}[1][]{%
  bfwbox,
  colback=bfwBgRose,
  colframe=bfwAccent!85!black,
  colbacktitle=bfwAccent!85!black,
  title={ESELSBRÜCKE},
  #1
}

% Beispiel-Box
\newtcolorbox{beispiel}[1][]{%
  bfwbox,
  colback=bfwBgSoft,
  colframe=bfwInkSoft,
  colbacktitle=bfwInkSoft,
  title={BEISPIEL},
  #1
}

% Lösungs-Box (für Aufgabenhefte)
\newtcolorbox{loesung}[1][]{%
  bfwbox,
  colback=bfwSuccess!7,
  colframe=bfwSuccess!65!black,
  colbacktitle=bfwSuccess!65!black,
  title={LÖSUNG},
  #1
}

% Achtung-Box — typische Fehler / Stolperfallen
\newtcolorbox{achtung}[1][]{%
  bfwbox,
  colback=bfwWarn!6,
  colframe=bfwWarn!85!black,
  colbacktitle=bfwWarn!85!black,
  title={ACHTUNG},
  #1
}

% Formel-Box — für zentrale Formeln (groß, ruhig, ohne Titelbalken)
\newtcolorbox{formelbox}[1][]{%
  enhanced,
  colback=bfwBgTeal,
  colframe=bfwPrimaryDark,
  boxrule=1pt, arc=3pt,
  left=10pt, right=10pt, top=6pt, bottom=6pt,
  #1
}

% Glossar-Eintrag
\newcommand{\glossareintrag}[2]{%
  \par\noindent\textbf{\color{bfwPrimaryDark}#1} \enspace #2\par\smallskip
}

% ---------------------------------------------------------------------------
% Abschnitts-Design: farbiges Nummern-Quadrat + feine Linie unter \section,
% dezent farbige \subsection/\subsubsection
% ---------------------------------------------------------------------------
\usepackage{titlesec}
\newcommand{\bfwSecBadge}[1]{%
  \begin{tikzpicture}[baseline={([yshift=-0.62em]n.center)}]
    \node[fill=bfwPrimaryDark, text=white, font=\normalsize\bfseries,
          minimum width=1.55em, minimum height=1.55em, inner sep=1pt,
          rounded corners=1.5pt] (n) {#1};
  \end{tikzpicture}}
\titleformat{\section}
  {\Large\bfseries\color{bfwPrimaryDark}}
  {\bfwSecBadge{\thesection}}{0.6em}{}
  [\vspace{-0.2em}{\color{bfwPrimary!60}\titlerule[0.8pt]}]
\titleformat{\subsection}
  {\large\bfseries\color{bfwPrimaryDark}}
  {\textcolor{bfwPrimary}{\thesubsection}}{0.5em}{}
\titleformat{\subsubsection}
  {\normalsize\bfseries\color{bfwInk}}
  {\textcolor{bfwPrimary}{\thesubsubsection}}{0.4em}{}
\titlespacing*{\section}{0pt}{1.6em}{1em}
\titlespacing*{\subsection}{0pt}{1.2em}{0.5em}

% TikZ-Stil "skizzenhaft" — etwas weicher als Rough.js, aber stabil
\tikzset{
  roughbox/.style = {
    draw=bfwPrimaryDark, thick, fill=bfwBgTeal,
    rounded corners=4pt, inner sep=6pt
  },
  roughaccent/.style = {
    draw=bfwAccent, thick, fill=bfwBgAmber,
    rounded corners=4pt, inner sep=6pt
  },
  roughline/.style = {
    draw=bfwInkSoft, thick, -{Stealth[length=2.5mm]}
  }
}

% Bullet-Symbole (bewusst kein FontAwesome — vermeidet Font-Installations-Probleme)
\newcommand{\faLightbulb}{$\bullet$}
\newcommand{\faBookmark}{$\bullet$}

% ---------------------------------------------------------------------------
% Kopf-/Fußzeile
%   Kopf links:  Wortlogo „Wirtschaftsfachwirt"
%   Kopf rechts: Dokument-Kurztext, gesetzt via \kopfzeile{...}
%   Fuß links:   © Can Siebert · 2026 — Fuß rechts: Seite X von Y
%   Titelseite (Seite 1) bleibt ohne Kopf-/Fußzeile (\handouttitel setzt
%   \thispagestyle{empty}).
% ---------------------------------------------------------------------------
\usepackage{fancyhdr}
\newcommand{\bfwKopfzeileText}{}
\newcommand{\kopfzeile}[1]{\renewcommand{\bfwKopfzeileText}{#1}}
\pagestyle{fancy}
\fancyhf{}
\fancyhead[L]{\small\bfseries\color{bfwPrimaryDark}Wirtschaftsfachwirt}
\fancyhead[R]{\small\color{bfwInkSoft}\bfwKopfzeileText}
\renewcommand{\headrulewidth}{0.6pt}
\renewcommand{\headrule}{{\color{bfwPrimary}\hrule width\headwidth height 0.6pt}}
\fancyfoot[L]{\small\color{bfwInkSoft}\copyright\ Can Siebert\,\textperiodcentered\,2026}
\fancyfoot[R]{\small\color{bfwInkSoft}Seite \thepage\ von \pageref*{LastPage}}
% Auch \thispagestyle{plain} (z.B. durch \maketitle) soll leer bleiben:
\fancypagestyle{plain}{\fancyhf{}\renewcommand{\headrulewidth}{0pt}\renewcommand{\headrule}{}}

% ---------------------------------------------------------------------------
% \handouttitel{Obertitel}{Untertitel}{Kurs-Label}{Termin-Zeile}
%   Einheitlicher professioneller Titelblock: Dachzeile, farbige Headline,
%   Untertitel, farbige Trennlinie und dezente Meta-Zeile (Kurs/Termin/Dozent).
%   Ersetzt \title/\maketitle. Direkt nach \begin{document} aufrufen.
% ---------------------------------------------------------------------------
\newcommand{\handouttitel}[4]{%
  \thispagestyle{empty}%
  \begingroup
  \setlength{\parindent}{0pt}%
  {\bfwlabelfont\color{bfwPrimary}WIRTSCHAFTSFACHWIRT (IHK)\par}
  \vspace{0.55em}
  {\huge\bfseries\color{bfwPrimaryDark}#1\par}
  \vspace{0.5em}
  {\large\color{bfwInkSoft}#2\par}
  \vspace{0.9em}
  {\color{bfwPrimary}\rule{\linewidth}{2pt}}\par
  \vspace{0.55em}
  {\small\color{bfwInkSoft}#3\ \,\textperiodcentered\,\ #4\ \,\textperiodcentered\,\ Dozent: Can Siebert\par}
  \endgroup
  \vspace{1.4em}%
}


\kopfzeile{Handout BWL Lektion 3 \textperiodcentered{} Teil 1}

\begin{document}
\handouttitel{Lektion~3 \textperiodcentered{} Teil~1: Existenzgründung \& Businessplan}%
  {Gründungsphasen, Voraussetzungen, Businessplan, Förderung und Kapitalbedarf}%
  {1.3 Existenzgründung \& Rechtsformen}%
  {Sa, 26.09.2026 \textperiodcentered{} Session 1 (08:15--10:40 Uhr)}

\begin{merksatz}[title={\bfseries Worum geht es in Teil 1?}]
Vor dem Schritt in die Selbstständigkeit sind viele Aufgaben zu bewältigen: die
Geschäftsidee ausarbeiten, Stärken und Schwächen analysieren, einen überzeugenden
Businessplan erstellen und die Finanzierung sichern. In dieser Session lernen Sie die
\emph{Gründungsphasen} und \emph{Voraussetzungen} der Existenzgründung kennen, erarbeiten
die \emph{Bestandteile des Businessplans}, verschaffen sich einen Überblick über die
\emph{Fördermöglichkeiten} (KfW, Gründungszuschuss) und berechnen den
\emph{Kapitalbedarf} einer Gründung. Die \emph{Unternehmensrechtsformen} folgen in
Teil~2 (Session~2, 10:55--13:15 Uhr). \emph{Hinweis: Der Auszug aus dem BWL-Textband
wird nachgereicht.}
\end{merksatz}

\tableofcontents
\vspace{1em}
\hrule
\vspace{1em}

\section{Gründungsphasen}

Die Gründung eines Unternehmens verläuft idealtypisch in mehreren Phasen. Am Anfang steht
die \textbf{Geschäftsidee} --- der Auslöser des gesamten Prozesses:

\begin{enumerate}
  \item \textbf{Geschäftsidee:} Eine Marktlücke, ein neues Produkt oder eine bessere
    Problemlösung wird erkannt. Die Idee muss ein Kundenproblem lösen und sich von
    bestehenden Angeboten abheben.
  \item \textbf{Orientierungsphase:} Erste Informationsbeschaffung und Selbstprüfung ---
    Passt die Gründung zu meiner Lebenssituation? Grobe Einschätzung von Markt, Wettbewerb
    und eigenem Können (u.\,a. Stärken-Schwächen-Analyse); Beratungsangebote (IHK,
    Gründungsberatung) nutzen.
  \item \textbf{Planungs- und Konzeptionsphase:} Aus der Idee wird ein Konzept --- der
    \emph{Businessplan} entsteht: Marktanalyse, Standort, Marketing, Rechtsformwahl,
    Kapitalbedarfs- und Finanzierungsplanung, Gespräche mit Kapitalgebern.
  \item \textbf{Umsetzungs- und Startphase:} Gründungsformalitäten (Gewerbeanmeldung,
    ggf.\ Handelsregistereintrag, Finanzamt, Kammern, Berufsgenossenschaft), Aufbau von
    Betrieb, Personal und Vertrieb, Markteintritt.
  \item \textbf{Festigungsphase (Konsolidierung):} Das junge Unternehmen etabliert sich am
    Markt --- Kundenstamm ausbauen, Liquidität sichern, Prozesse professionalisieren,
    Soll-Ist-Kontrolle gegen den Businessplan, ggf.\ nachsteuern und wachsen.
\end{enumerate}

\begin{eselsbruecke}
\textbf{G-O-K-U-F}: \textbf{G}eschäftsidee $\rightarrow$ \textbf{O}rientierung
$\rightarrow$ \textbf{K}onzeption $\rightarrow$ \textbf{U}msetzung $\rightarrow$
\textbf{F}estigung. („Gute Organisation kann Unternehmen festigen.``)
\end{eselsbruecke}

\begin{beispiel}
Das Gründungsgeschehen in Deutschland schwankt mit der Konjunktur und dem Arbeitsmarkt:
Der \emph{Hochpunkt 2004} bei den Existenzgründungen wird u.\,a. mit den
arbeitsmarktpolitischen Anreizen der „Hartz-Gesetze`` erklärt (Gründung aus der
Arbeitslosigkeit), der \emph{Tiefpunkt 2015} mit der guten Arbeitsmarktlage, dem
Fachkräftemangel und der demografischen Entwicklung --- wer eine sichere Anstellung
findet, gründet seltener (vgl.\ Übungsband, Übung~45).
\end{beispiel}

\section{Voraussetzungen der Existenzgründung}

Nicht jede Person und nicht jede Idee eignet sich für eine Gründung. Vor dem Start sind
vier Bereiche kritisch zu prüfen:

\subsection{Persönliche Ausgangslage}

\begin{itemize}
  \item Fähigkeit, sich selbst Ziele zu setzen und diese konsequent zu verfolgen;
  \item Mut zum Risiko und Fähigkeit, mit Rückschlägen fertig zu werden (Belastbarkeit);
  \item Rückhalt der Familie und Bereitschaft zu hohem zeitlichen Einsatz;
  \item Disziplin im Umgang mit Geld und Bereitschaft zu finanziellen Einschränkungen
    in der Anfangsphase; Gesundheit und Eigenmotivation.
\end{itemize}

\subsection{Fachliche und kaufmännische Qualifikation}

\begin{itemize}
  \item \textbf{Fachliche Kompetenz:} Ausbildung und Berufserfahrung passen zum
    Gründungsvorhaben; Branchenkenntnis (Entwicklung, Kunden, Wettbewerber); Bereitschaft,
    fachliche Defizite auszugleichen; Umgang mit branchenspezifischer EDV.
  \item \textbf{Kaufmännisch-unternehmerische Kompetenz:} Kenntnisse in Kalkulation,
    Buchführung, Steuern, Recht und Vertrieb; selbstbewusstes und kontaktfreudiges
    Auftreten; Fähigkeit zu delegieren und anderen zu vertrauen; Kreativität, um das
    Unternehmen voranzubringen.
\end{itemize}

\subsection{Finanzielle Basis}

Ausreichendes \textbf{Eigenkapital} (Faustregel: mindestens 15--20\,\% des
Kapitalbedarfs), realistische Einschätzung des \textbf{Kapitalbedarfs} (Investitionen,
Anlaufkosten, private Lebenshaltung in der Startphase) sowie Zugang zu
\textbf{Fremdkapital} und Sicherheiten. Häufigste Ursache des Scheiterns junger
Unternehmen sind Finanzierungsmängel und unzureichende Planung.

\subsection{Tragfähige Geschäftsidee}

Die Idee muss einen \emph{Kundennutzen} stiften, sich vom Wettbewerb \emph{abgrenzen}
(Alleinstellungsmerkmal, USP) und auf einem ausreichend großen, erreichbaren \emph{Markt}
umsetzbar sein.

\begin{merksatz}
Die IHK-Prüfung fragt die Voraussetzungen gern in drei Gruppen ab: \textbf{persönliche
Voraussetzungen}, \textbf{fachliche Kompetenz} und \textbf{unternehmerische
(kaufmännische) Kompetenz} --- nennen Sie je drei bis vier konkrete Punkte
(vgl.\ Übungsband, Übung~49).
\end{merksatz}

\section{Der Businessplan}

\subsection{Begriff, Funktionen und Adressaten}

Der \textbf{Businessplan} (Geschäftsplan) ist die \emph{schriftliche} Ausarbeitung zur
Darlegung der \emph{Geschäftsidee}, ihrer Realisierung und der daraus resultierenden
Ergebnisse. Er erfüllt drei Funktionen:

\begin{itemize}
  \item \textbf{Nach außen:} Er richtet sich v.\,a.\ an die \emph{Geldgeber} (Banken,
    Förderinstitute, Beteiligungsgeber) und zukünftige Geschäftspartner --- ohne
    überzeugenden Businessplan keine Finanzierung.
  \item \textbf{Nach innen (Planung):} Er zwingt den Existenzgründer, die Geschäftsidee
    \emph{planmäßig} zu durchdenken --- Schwachstellen werden früh sichtbar.
  \item \textbf{Nach innen (Steuerung):} Er dient als \emph{Führungsgrundlage} für
    Mitarbeiter und Gremien im zu gründenden Unternehmen und später als Maßstab für
    Soll-Ist-Vergleiche.
\end{itemize}

\subsection{Bestandteile des Businessplans}

\begin{enumerate}
  \item \textbf{Executive Summary:} Kurzzusammenfassung des gesamten Plans auf 1--2
    Seiten --- wird zuerst gelesen und entscheidet oft über das Weiterlesen.
  \item \textbf{Geschäftsmodell/Geschäftsidee:} Zweck des Vorhabens, Besonderheit der
    Geschäftsidee, Art des Angebots (Produkt/Dienstleistung), Kundennutzen,
    Unternehmensziele; Gründerprofil (Qualifikation des Gründungsteams).
  \item \textbf{Markt und Wettbewerb:} Markt- und Branchenanalyse (Marktgröße, -wachstum,
    Zielgruppen), Wettbewerbsanalyse (Stärken/Schwächen der Konkurrenz), Standortanalyse.
  \item \textbf{Marketing und Vertrieb:} Marketing-Mix --- Produkt-, Preis-,
    Distributions- und Kommunikationspolitik; geplante Absatzwege.
  \item \textbf{Organisation und Rechtsform:} Aufbauorganisation, Personalplanung,
    gewählte Rechtsform mit Begründung (dazu ausführlich Teil~2 dieser Lektion).
  \item \textbf{Finanzplanung und Kapitalbedarf:} Kalkulation des \emph{Kapitalbedarfs},
    \emph{Investitionsplan}, \emph{Finanzierungsplan} (Eigen-/Fremdkapital, Fördermittel),
    \emph{Liquiditätsplan}, \emph{Umsatzschätzung und Ertragsvorschau}
    (Rentabilitätsplanung).
  \item \textbf{Chancen und Risiken:} realistische Szenarien (Best/Worst Case) und
    Gegenmaßnahmen.
\end{enumerate}

\begin{eselsbruecke}
Die Kernbestandteile als Kette: \textbf{„Erst GeMMOF-C prüfen!``} ---
\textbf{E}xecutive Summary, \textbf{Ge}schäftsmodell, \textbf{M}arkt/Wettbewerb,
\textbf{M}arketing, \textbf{O}rganisation/Rechtsform, \textbf{F}inanzplanung,
\textbf{C}hancen/Risiken.
\end{eselsbruecke}

\subsection{Fördermöglichkeiten (Überblick)}

\begin{itemize}
  \item \textbf{KfW-Förderung:} zinsgünstige Gründerkredite der staatlichen KfW-Bank
    (z.\,B. ERP-Gründerkredit „StartGeld`` bis 125.000~€), beantragt über die Hausbank
    (Hausbankprinzip); die KfW übernimmt dabei einen großen Teil des Kreditrisikos.
  \item \textbf{Gründungszuschuss:} Zuschuss der \emph{Agentur für Arbeit} für Gründungen
    aus dem Bezug von Arbeitslosengeld~I (Ermessensleistung; ALG-I-Betrag plus Pauschale
    zur sozialen Absicherung, Tragfähigkeitsbescheinigung z.\,B. der IHK erforderlich).
  \item Daneben: Bürgschaftsbanken (Sicherheiten), Landesförderprogramme,
    Beratungsförderung und Beteiligungskapital.
\end{itemize}

\section{Formel \& Rechenbeispiel: Kapitalbedarfsrechnung}

Der Kapitalbedarf wird als einfaches \emph{Additionsschema} ermittelt --- erst danach
zeigt sich, wie viel Fremdkapital fehlt:

\[
\textrm{Kapitalbedarf} = \textrm{Investitionen} + \textrm{Gründungskosten}
+ \textrm{Anlaufkosten} + \textrm{private Lebenshaltung}
\]
\[
\textrm{Fremdkapitalbedarf} = \textrm{Kapitalbedarf} - \textrm{Eigenkapital}
\]

\begin{beispiel}[title={\bfseries Zahlenbeispiel: Kapitalbedarf einer Ladengründung}]
\begin{center}
\begin{tabular}{clr}
\toprule
 & \textbf{Position} & \textbf{Betrag}\\
\midrule
 & Investitionen (Ladeneinrichtung, Maschinen) & 60.000~€\\
$+$ & Warenerstausstattung & 20.000~€\\
$+$ & Gründungskosten (Notar, Register, Beratung) & 3.000~€\\
$+$ & Anlaufkosten (6 Monate Miete, Personal, Werbung) & 15.000~€\\
$+$ & private Lebenshaltung (6 Monate) & 12.000~€\\
\midrule
$=$ & \textbf{Kapitalbedarf} & \textbf{110.000~€}\\
$-$ & Eigenkapital & 22.000~€\\
\midrule
$=$ & \textbf{Fremdkapitalbedarf} & \textbf{88.000~€}\\
\bottomrule
\end{tabular}
\end{center}
\textbf{Rechenweg:} Schritt~1: alle Positionen addieren (110.000~€).
Schritt~2: Eigenkapital abziehen (110.000 $-$ 22.000 = 88.000~€ Fremdkapitalbedarf).
Schritt~3: Eigenkapitalquote prüfen: $22.000 \div 110.000 = 20\,\%$ ---
die Faustregel (15--20\,\%) ist erfüllt.
\end{beispiel}

\begin{merksatz}[title={\bfseries Wichtig für die Prüfungsvorbereitung (Teil 1)}]
Sie sollten sicher können: die \emph{fünf Gründungsphasen} in der richtigen Reihenfolge,
die \emph{Voraussetzungen} der Gründung in drei Gruppen, die \emph{drei Funktionen} und
die \emph{Bestandteile des Businessplans} mit den Elementen der Finanzplanung, die
\emph{Förderwege} (KfW über die Hausbank, Gründungszuschuss der Agentur für Arbeit)
sowie die \emph{Kapitalbedarfsrechnung}. In Teil~2 folgen die Unternehmensrechtsformen
--- dort wird auch die Rechtsformwahl für den Businessplan-Baustein „Organisation und
Rechtsform`` begründet.
\end{merksatz}

\vspace{1em}
\hrule
\vspace{0.8em}
\noindent\textsf{\small Quellen: Rahmenplan Wirtschaftsbezogene Qualifikationen (DIHK), Abschnitt 1.3;
Übungsband Betriebswirtschaft (DIHK-Bildungs-GmbH), Kap.~2, S.~21--24 (Lösungen S.~45--47);
Textband-Auszug wird nachgereicht. Weiter geht es in Teil~2 (Session~2, 10:55--13:15 Uhr):
Unternehmensrechtsformen.}

\end{document}
